% Corrected French for printed p.140.
% Mathematical correction: the universal extension is from degree zero.
% Corrections: remove duplicated phrase in the naturality condition;
% C -> C' as codomain of F; foncteur -> foncteurs before universels.
\SourcePage{140}

\LayerTarget{tohoku:axiom:delta-functor:2:p2}
\noindent
\[
\tag{2.1.1.}
\ldots\,.\;T^i(A')\to T^i(A)
\to T^i(A'')\to T^{i+1}(A')
\to\ldots\,.
\]

\noindent est un complexe, i. e. le composé de deux morphismes consécutifs
dans cette suite est nul.

\LayerTarget{tohoku:def:delta-star-and-variance}
Définition analogue pour un $\partial^*$-foncteur covariant, la seule
différence étant que l'opérateur $\partial$ diminue le degré d'une unité au
lieu de l'augmenter. Définition analogue pour les $\partial$-foncteurs et
$\partial^*$-foncteurs contravariants ; les $T^i$ sont alors des foncteurs
additifs contravariants et les opérateurs bords vont de $T^i(A')$ dans
$T^{i+1}(A'')$ ou $T^{i-1}(A'')$. Si on change de signe les degrés $i$ dans
$T^i$, ou si on remplace la catégorie $\mathbf{C}'$ par sa duale, les
$\partial$-foncteurs sont changés en $\partial^*$-foncteurs. Ainsi, on peut
toujours se ramener à l'étude des $\partial$-foncteurs covariants. Remarquons
que si $a=-\infty$, $b=+\infty$, un $\partial$-foncteur est une
\textquotedblleft connected sequence of functors\textquotedblright\ de [6,
Chap. III].

\LayerTarget{tohoku:def:delta-functor-morphism}
Etant donnés deux $\partial$-foncteurs $T$ et $T'$ définis pour les mêmes
degrés, on appelle \textit{morphisme} (ou transformation naturelle) de $T$
dans $T'$ un système $f=(f^i)$ de morphismes fonctoriels
$f^i:T^i\to T'^i$, soumis à la condition naturelle de commutation
avec $\partial$ : pour toute suite exacte
$0\to A'\to A\to A''\to 0$,
le diagramme

\LayerTarget{tohoku:diagram:delta-morphism}
\[
\begin{array}{ccc}
T^i(A'') & \mathop{\to}\limits^{\partial} & T^{i+1}(A')\\
\big\downarrow && \big\downarrow\\
T'^i(A'') & \mathop{\to}\limits^{\partial} & T'^{i+1}(A')
\end{array}
\]

\noindent est commutatif. Les morphismes de $\partial$-foncteurs s'ajoutent
et se composent de façon évidente.

\LayerTarget{tohoku:def:exact-delta-functor}
Supposons que $\mathbf{C}'$ soit aussi une catégorie abélienne. Alors un
$\partial$-foncteur est dit \textit{exact} si pour toute suite exacte
$0\to A'\to A\to A''\to 0$
dans $\mathbf{C}$, la suite correspondante (2.1.1) de morphismes est exacte. On appelle
\textit{foncteur cohomologique} (resp. \textit{foncteur homologique}) un
$\partial$-foncteur (resp. $\partial^*$-foncteur) \textit{exact} défini pour tous les
degrés.

\LayerTarget{tohoku:section:2.2}
\indent\textbf{2.2. $\partial$-foncteurs universels.}
\LayerTarget{tohoku:def:universal-delta-functor}
Soit $T=(T^i)\quad 0\leqq i\leqq a$ un $\partial$-foncteur covariant de $\mathbf{C}$ dans
$\mathbf{C}'$, où $a>0$. $T$ est dit un $\partial$-foncteur \textit{universel} si pour
tout $\partial$-foncteur $T'=(T'^i)$ défini pour les mêmes degrés, et tout
morphisme fonctoriel $f^0$ de $T^0$ dans $T'^0$, il existe un
morphisme et un seul $f$ de $T$ dans $T'$ qui pour le degré $0$ se réduise à
$f^0$. Définition identique pour un $\partial$-foncteur contravariant ;
dans le cas des $\partial^*$-foncteurs il faut considérer des morphismes de
$T'$ dans $T$ et non de $T$ dans $T'$.

\LayerTarget{tohoku:prop:satellites-uniqueness}
Par définition même, pour un foncteur covariant $F$ donné de $\mathbf{C}$ dans $\mathbf{C}'$, et
pour $a$, donné, $a>0$, il ne peut exister, à isomorphismes canoniques près,
qu'un seul $\partial$-foncteur universel défini en degrés $0\leqq i\leqq a$ et
se réduisant à $F$ en degré 0.
\LayerTarget{tohoku:def:right-left-satellites}
Ses composantes sont alors notées $S^iF$, et appelées \textit{foncteurs
satellites} droits de $F$. Si $i>0$, on posera aussi $S_iF=S^{-i}F$, où les
$S_iF$ sont les satellites \textit{gauches} de $F$, définis comme $S^iF$,
mais par la considération des $\partial^*$-foncteurs universels définis en
degrés $0\leqq i\leqq a$, tels que $T^0=F$. On vérifie aussitôt que pour $i$
donné, si $S^iF$ existe, il est indépendant du choix de $a$.

\LayerTarget{tohoku:remark:satellites-existence:p1}
Dans tous les cas à ma connaissance, les foncteurs satellites d'un foncteur
additif quelconque $F$ existent. D'ailleurs, si $\mathbf{C},\mathbf{C}'$ sont donnés, pour
montrer

\TohokuContinuation{141}
